%% Template beamer nao oficial com a identidade visual da FGV.
%% Baseado no template do IMPA de Rafael Vieira Westenberger (CC BY-SA 4.0),
%% totalmente reescrito segundo o "Manual de Identidade Visual - Conveniadas
%% IDE" (FGV/SCM, maio de 2009):
%%   * Cor padrao: azul escuro, Pantone 281C, RGB 28/47/103 (pag. 11)
%%   * Tipografia: Arial e derivados para documentos (pag. 10); aqui,
%%     Helvetica (helvet) em pdfLaTeX; titulos em versao condensada,
%%     como no manual (Zurich Cn BT)
%%   * Marca FGV (simbolo V + "FGV"): vetor extraido do proprio manual,
%%     embutido abaixo em \fgvlogo -- nao e preciso nenhum arquivo externo
%%   * Layout inspirado no manual: faixa de titulo com filete azul,
%%     rodape com filete duplo, capa com faixa lateral azul, caixas
%%     lavanda com seta
%% Compile com pdfLaTeX (ou latexmk -pdf).

\documentclass[aspectratio=1610,10pt]{beamer}
\usepackage[T1]{fontenc}
\usepackage[utf8]{inputenc}
\usepackage[portuguese]{babel} % Garante a hifenização correta e tradução de termos
\usepackage{amsmath, amssymb, amsthm, mathtools} % Rigor matemático e símbolos
\usepackage{hyperref} % Links interativos
\usepackage{helvet}
\renewcommand{\familydefault}{\sfdefault}
\usepackage{microtype}
\usepackage{tikz}
\usefonttheme[onlymath]{serif}

% Configuração para sumário e links interativos discretos
\hypersetup{
    colorlinks=true,
    linkcolor=main,
    urlcolor=main,
    citecolor=main
}

% Definição de ambientes matemáticos formais em português
\newtheorem{teorema}{Teorema}
\newtheorem{lema}{Lema}
\newtheorem{definicao}{Definição}

%% ------------------------------------------------------------------
%% Cores (manual FGV)
%% ------------------------------------------------------------------
\definecolor{main}{RGB}{28,47,103}        % azul padrao FGV (Pantone 281C)
\definecolor{fgvcyan}{RGB}{0,173,239}     % numerais/destaques dos titulos
\definecolor{fgvlav}{RGB}{192,194,223}    % barras lavanda dos subtitulos
\definecolor{fgvgray}{RGB}{120,126,138}   % notas laterais / numeracao
\definecolor{fgvink}{RGB}{35,31,32}       % texto corrente

%% ------------------------------------------------------------------
%% Marca FGV (vetor do manual, pag. 3). \fgvlogo[cor]{escala}
%% ------------------------------------------------------------------
\newcommand{\fgvlogo}[2][main]{%
  \begin{pgfpicture}%
    \pgftransformscale{#2}%
    \pgfsetfillcolor{#1}%
\pgfpathmoveto{\pgfpoint{20.266pt}{22.020pt}}%
\pgfpathlineto{\pgfpoint{26.312pt}{26.008pt}}%
\pgfpathlineto{\pgfpoint{37.504pt}{26.008pt}}%
\pgfpathlineto{\pgfpoint{25.953pt}{18.344pt}}%
\pgfpathclose%
\pgfpathmoveto{\pgfpoint{0.000pt}{26.008pt}}%
\pgfpathlineto{\pgfpoint{10.996pt}{26.008pt}}%
\pgfpathlineto{\pgfpoint{22.977pt}{12.234pt}}%
\pgfpathlineto{\pgfpoint{11.980pt}{12.234pt}}%
\pgfpathclose%
\pgfpathmoveto{\pgfpoint{9.211pt}{0.141pt}}%
\pgfpathlineto{\pgfpoint{9.211pt}{9.203pt}}%
\pgfpathlineto{\pgfpoint{13.883pt}{9.203pt}}%
\pgfpathlineto{\pgfpoint{13.883pt}{7.805pt}}%
\pgfpathlineto{\pgfpoint{11.121pt}{7.805pt}}%
\pgfpathlineto{\pgfpoint{11.121pt}{5.641pt}}%
\pgfpathlineto{\pgfpoint{13.754pt}{5.641pt}}%
\pgfpathlineto{\pgfpoint{13.754pt}{4.246pt}}%
\pgfpathlineto{\pgfpoint{11.121pt}{4.246pt}}%
\pgfpathlineto{\pgfpoint{11.121pt}{0.141pt}}%
\pgfpathclose%
\pgfpathmoveto{\pgfpoint{20.961pt}{4.910pt}}%
\pgfpathlineto{\pgfpoint{20.961pt}{0.309pt}}%
\pgfpathcurveto{\pgfpoint{20.352pt}{0.199pt}}{\pgfpoint{19.836pt}{0.117pt}}{\pgfpoint{19.410pt}{0.070pt}}%
\pgfpathcurveto{\pgfpoint{18.988pt}{0.023pt}}{\pgfpoint{18.605pt}{0.000pt}}{\pgfpoint{18.262pt}{0.000pt}}%
\pgfpathcurveto{\pgfpoint{17.707pt}{0.000pt}}{\pgfpoint{17.242pt}{0.051pt}}{\pgfpoint{16.859pt}{0.152pt}}%
\pgfpathcurveto{\pgfpoint{16.477pt}{0.258pt}}{\pgfpoint{16.164pt}{0.410pt}}{\pgfpoint{15.918pt}{0.621pt}}%
\pgfpathcurveto{\pgfpoint{15.551pt}{0.949pt}}{\pgfpoint{15.285pt}{1.410pt}}{\pgfpoint{15.129pt}{2.004pt}}%
\pgfpathcurveto{\pgfpoint{14.973pt}{2.602pt}}{\pgfpoint{14.895pt}{3.562pt}}{\pgfpoint{14.895pt}{4.887pt}}%
\pgfpathcurveto{\pgfpoint{14.895pt}{5.797pt}}{\pgfpoint{14.941pt}{6.508pt}}{\pgfpoint{15.039pt}{7.012pt}}%
\pgfpathcurveto{\pgfpoint{15.137pt}{7.516pt}}{\pgfpoint{15.297pt}{7.934pt}}{\pgfpoint{15.520pt}{8.270pt}}%
\pgfpathcurveto{\pgfpoint{15.754pt}{8.617pt}}{\pgfpoint{16.086pt}{8.879pt}}{\pgfpoint{16.508pt}{9.055pt}}%
\pgfpathcurveto{\pgfpoint{16.930pt}{9.230pt}}{\pgfpoint{17.449pt}{9.316pt}}{\pgfpoint{18.066pt}{9.316pt}}%
\pgfpathcurveto{\pgfpoint{19.008pt}{9.316pt}}{\pgfpoint{19.727pt}{9.078pt}}{\pgfpoint{20.219pt}{8.609pt}}%
\pgfpathcurveto{\pgfpoint{20.711pt}{8.137pt}}{\pgfpoint{20.961pt}{7.449pt}}{\pgfpoint{20.961pt}{6.543pt}}%
\pgfpathlineto{\pgfpoint{20.961pt}{6.406pt}}%
\pgfpathlineto{\pgfpoint{19.102pt}{6.406pt}}%
\pgfpathlineto{\pgfpoint{19.102pt}{6.520pt}}%
\pgfpathcurveto{\pgfpoint{19.102pt}{7.031pt}}{\pgfpoint{19.008pt}{7.410pt}}{\pgfpoint{18.832pt}{7.660pt}}%
\pgfpathcurveto{\pgfpoint{18.652pt}{7.906pt}}{\pgfpoint{18.379pt}{8.031pt}}{\pgfpoint{18.012pt}{8.031pt}}%
\pgfpathcurveto{\pgfpoint{17.547pt}{8.031pt}}{\pgfpoint{17.230pt}{7.824pt}}{\pgfpoint{17.055pt}{7.414pt}}%
\pgfpathcurveto{\pgfpoint{16.883pt}{7.000pt}}{\pgfpoint{16.793pt}{6.094pt}}{\pgfpoint{16.793pt}{4.695pt}}%
\pgfpathcurveto{\pgfpoint{16.793pt}{3.215pt}}{\pgfpoint{16.898pt}{2.285pt}}{\pgfpoint{17.109pt}{1.902pt}}%
\pgfpathcurveto{\pgfpoint{17.320pt}{1.520pt}}{\pgfpoint{17.715pt}{1.328pt}}{\pgfpoint{18.293pt}{1.328pt}}%
\pgfpathcurveto{\pgfpoint{18.422pt}{1.328pt}}{\pgfpoint{18.559pt}{1.340pt}}{\pgfpoint{18.707pt}{1.359pt}}%
\pgfpathcurveto{\pgfpoint{18.859pt}{1.379pt}}{\pgfpoint{19.031pt}{1.414pt}}{\pgfpoint{19.234pt}{1.461pt}}%
\pgfpathlineto{\pgfpoint{19.234pt}{3.488pt}}%
\pgfpathlineto{\pgfpoint{18.191pt}{3.488pt}}%
\pgfpathlineto{\pgfpoint{18.191pt}{4.910pt}}%
\pgfpathclose%
\pgfpathmoveto{\pgfpoint{23.707pt}{0.129pt}}%
\pgfpathlineto{\pgfpoint{21.430pt}{9.203pt}}%
\pgfpathlineto{\pgfpoint{23.406pt}{9.203pt}}%
\pgfpathlineto{\pgfpoint{24.336pt}{4.473pt}}%
\pgfpathcurveto{\pgfpoint{24.488pt}{3.703pt}}{\pgfpoint{24.605pt}{3.102pt}}{\pgfpoint{24.680pt}{2.664pt}}%
\pgfpathcurveto{\pgfpoint{24.750pt}{2.230pt}}{\pgfpoint{24.805pt}{1.840pt}}{\pgfpoint{24.836pt}{1.500pt}}%
\pgfpathcurveto{\pgfpoint{24.898pt}{2.059pt}}{\pgfpoint{24.969pt}{2.590pt}}{\pgfpoint{25.051pt}{3.090pt}}%
\pgfpathcurveto{\pgfpoint{25.129pt}{3.586pt}}{\pgfpoint{25.219pt}{4.047pt}}{\pgfpoint{25.316pt}{4.473pt}}%
\pgfpathlineto{\pgfpoint{26.398pt}{9.203pt}}%
\pgfpathlineto{\pgfpoint{28.309pt}{9.203pt}}%
\pgfpathlineto{\pgfpoint{25.941pt}{0.129pt}}%
\pgfpathlineto{\pgfpoint{23.707pt}{0.129pt}}%
    \pgfseteorule\pgfusepath{fill}%
  \end{pgfpicture}}

%% ------------------------------------------------------------------
%% Tipografia
%% ------------------------------------------------------------------
\newcommand{\fgvcondensed}{\fontfamily{phv}\fontshape{n}\selectfont}

\setbeamerfont{title}{size=\huge,series=\bfseries}
\setbeamerfont{subtitle}{size=\large,series=\mdseries}
\setbeamerfont{frametitle}{size=\LARGE,series=\mdseries}
\setbeamerfont{block title}{size=\normalsize,series=\bfseries}
\setbeamerfont{footline}{size=\scriptsize}

%% ------------------------------------------------------------------
%% Cores do beamer
%% ------------------------------------------------------------------
\setbeamercolor{normal text}{fg=fgvink,bg=white}
\setbeamercolor{structure}{fg=main}
\setbeamercolor{alerted text}{fg=fgvcyan!85!black}
\setbeamercolor{example text}{fg=main}
\setbeamercolor{frametitle}{fg=main}
\setbeamercolor{title}{fg=fgvink}
\setbeamercolor{subtitle}{fg=main}
\setbeamercolor{author}{fg=fgvink}
\setbeamercolor{institute}{fg=fgvgray}
\setbeamercolor{date}{fg=fgvgray}
\setbeamercolor{block title}{fg=main,bg=fgvlav}
\setbeamercolor{block body}{fg=fgvink,bg=fgvlav!22}
\setbeamercolor{block title alerted}{fg=main,bg=fgvcyan!35}
\setbeamercolor{block body alerted}{fg=fgvink,bg=fgvcyan!8}
\setbeamercolor{block title example}{fg=main,bg=fgvlav}
\setbeamercolor{block body example}{fg=fgvink,bg=fgvlav!22}
\setbeamercolor{item}{fg=main}
\setbeamercolor{caption name}{fg=main}
\setbeamercolor{footnote}{fg=fgvgray}
\setbeamercolor{section in toc}{fg=main}
\setbeamercolor{section in toc shaded}{fg=main!55}
\setbeamercolor{subsection in toc shaded}{fg=fgvink!55}
\setbeamercolor{subsection in toc}{fg=fgvink}
\setbeamercolor{section number projected}{bg=main,fg=white}

%% ------------------------------------------------------------------
%% Marcadores
%% ------------------------------------------------------------------
\setbeamertemplate{navigation symbols}{}
\setbeamertemplate{itemize item}{\scriptsize$\blacktriangleright$}
\setbeamertemplate{itemize subitem}{\tiny$\blacktriangleright$}
\setbeamertemplate{itemize subsubitem}{\tiny$\blacktriangleright$}
\setbeamertemplate{enumerate item}{\insertenumlabel.}
\setbeamertemplate{enumerate subitem}{\insertenumlabel.\insertsubenumlabel}
\setbeamertemplate{caption}[numbered]
\setbeamertemplate{caption label separator}{: }
\setbeamertemplate{frametitle continuation}{}

%% ------------------------------------------------------------------
%% Margens
%% ------------------------------------------------------------------
\setbeamersize{text margin left=1.4cm,text margin right=1.4cm}

%% ------------------------------------------------------------------
%% Barra de titulo
%% ------------------------------------------------------------------
\newcommand{\fgv}{}   
\newif\iffgvnumber    
\fgvnumberfalse

\setbeamertemplate{headline}{%
  \ifnum\value{page}=1\relax\else
  \leavevmode
  \begin{tikzpicture}[remember picture,overlay]
    \shade[top color=fgvlav!45,bottom color=white]
      (current page.north west) rectangle ([yshift=-1.55cm]current page.north east);
    \fill[main] ([yshift=-1.55cm]current page.north west)
      rectangle ([yshift=-1.63cm,xshift=-2.35cm]current page.north east);
    \node[anchor=east,inner sep=0pt] at
      ([xshift=-1.0cm,yshift=-0.78cm]current page.north east)
      {\fgvlogo{0.62}};
  \end{tikzpicture}%
  \fi}

\setbeamertemplate{frametitle}{%
  \vbox{%
    \vskip0.30cm
    \begin{beamercolorbox}[wd=\paperwidth,ht=1.0cm,dp=0.2cm,leftskip=1.4cm,rightskip=3.4cm]{frametitle}%
      \fgvcondensed\usebeamerfont{frametitle}\usebeamercolor[fg]{frametitle}%
      \iffgvnumber\ifnum\value{section}>0
        {\color{fgvcyan}\fontseries{b}\selectfont\thesection.}\,%
      \fi\fi
      \insertframetitle\par
    \end{beamercolorbox}%
    \vskip0.40cm}}

%% ------------------------------------------------------------------
%% Rodape
%% ------------------------------------------------------------------
\setbeamertemplate{footline}{%
  \leavevmode
  \begin{tikzpicture}[remember picture,overlay]
    \fill[fgvlav]
      ([yshift=0.50cm]current page.south west) rectangle
      ([yshift=0.455cm]current page.south east);
    \fill[main]
      ([yshift=0.60cm]current page.south west) rectangle
      ([yshift=0.54cm,xshift=-0.3cm]current page.south east);
    \node[anchor=east,inner sep=0pt,text=fgvgray] at
      ([xshift=-0.6cm,yshift=0.27cm]current page.south east)
      {\usebeamerfont{footline}\insertframenumber};
    \node[anchor=west,inner sep=0pt,text=fgvgray] at
      ([xshift=1.4cm,yshift=0.27cm]current page.south west)
      {\usebeamerfont{footline}\insertshortauthor};
  \end{tikzpicture}%
  \vskip0.7cm}

%% ------------------------------------------------------------------
%% Blocos
%% ------------------------------------------------------------------
\setbeamertemplate{blocks}[default]
\setbeamertemplate{block begin}{%
  \par\vskip\medskipamount
  \begin{beamercolorbox}[colsep*=0pt,leftskip=0pt,rightskip=0pt,wd=\textwidth]{block title}%
    \begin{tikzpicture}[baseline=(t.base)]
      \node[inner sep=0pt] (t) {\strut};
    \end{tikzpicture}%
    \hspace*{0.22cm}\raisebox{0.02cm}{\color{main}$\Rightarrow$}\hspace*{0.2cm}%
    \usebeamerfont{block title}\insertblocktitle\strut\hfill
  \end{beamercolorbox}%
  \nointerlineskip
  \begin{beamercolorbox}[colsep=0.25cm,vmode,wd=\textwidth]{block body}%
    \ifbeamercolorempty[bg]{block body}{}{\vskip0.05cm}
    \usebeamerfont{block body}}
\setbeamertemplate{block end}{%
    \vskip0.05cm
  \end{beamercolorbox}%
  \vskip\medskipamount}

%% ------------------------------------------------------------------
%% Pagina de titulo
%% ------------------------------------------------------------------
\newcommand{\fgvband}{Apresentação}   % texto vertical da faixa

\setbeamertemplate{title page}{%
  \begin{tikzpicture}[remember picture,overlay]
    \shade[top color=fgvlav!35,bottom color=white,shading angle=0]
      (current page.south west) rectangle (current page.north east);
    \begin{scope}[opacity=0.07]
      \node[inner sep=0pt,anchor=center] at
        ([xshift=1.0cm,yshift=3.3cm]current page.center) {\fgvlogo[main]{4.2}};
    \end{scope}
    \fill[main] ([xshift=-3.3cm]current page.north east) rectangle
      (current page.south east);
    \node[rotate=90,text=white,inner sep=0pt,anchor=west] at
      ([xshift=-1.65cm,yshift=3.7cm]current page.south east)
      {\fgvcondensed\itshape\fontsize{20}{20}\selectfont\MakeUppercase{\fgvband}};
    \node[anchor=south,inner sep=0pt] at
      ([xshift=-1.65cm,yshift=1.1cm]current page.south east)
      {\fgvlogo[white]{1.35}};
    \node[anchor=south west,text width=\paperwidth-7.4cm,inner sep=0pt] at
      ([xshift=1.4cm,yshift=5.25cm]current page.south west)
      {\raggedright\usebeamercolor[fg]{title}\fgvcondensed\usebeamerfont{title}\inserttitle\par
       \vspace{0.15cm}
       \usebeamercolor[fg]{subtitle}\usebeamerfont{subtitle}\insertsubtitle};
    \fill[main] ([xshift=1.4cm,yshift=4.95cm]current page.south west)
      rectangle ([xshift=-3.3cm,yshift=4.88cm]current page.south east);
    \node[anchor=north west,text width=\paperwidth-7.4cm,inner sep=0pt] at
      ([xshift=1.4cm,yshift=4.65cm]current page.south west)
      {\raggedright\usebeamercolor[fg]{author}\insertauthor\par
       \vspace{0.2cm}
       \usebeamercolor[fg]{institute}\footnotesize\insertinstitute\par
       \vspace{0.15cm}
       \usebeamercolor[fg]{date}\insertdate};
  \end{tikzpicture}}

%% ------------------------------------------------------------------
%% Quadro de destaque
%% ------------------------------------------------------------------
\newenvironment{fgvfocus}{%
  \begin{tikzpicture}[remember picture,overlay]
    \fill[main] (current page.south west) rectangle (current page.north east);
    \node[anchor=south east,inner sep=0pt] at
      ([xshift=-1.0cm,yshift=0.8cm]current page.south east)
      {\fgvlogo[white]{1.1}};
  \end{tikzpicture}%
  \color{white}%
  \begin{center}\fgvcondensed\fontsize{30}{36}\selectfont
}{\end{center}}

%% ------------------------------------------------------------------
%% Sumario por secao
%% ------------------------------------------------------------------
\setbeamertemplate{section in toc}{%
  {\fgvcondensed\bfseries\color{main}\MakeUppercase{Capítulo}\ \inserttocsectionnumber\ ::\ \inserttocsection}\par}
\setbeamertemplate{subsection in toc}{%
  \hspace*{1.2em}\inserttocsectionnumber.\inserttocsubsectionnumber\hspace{0.8em}\inserttocsubsection\par}

\AtBeginSection[]
{
  \begin{frame}<beamer>
    \frametitle{Sumário}
    \tableofcontents[currentsection,hideothersubsections]
  \end{frame}
}
\AtBeginSubsection[]{}

%% ------------------------------------------------------------------
%% Dados da apresentacao
%% ------------------------------------------------------------------
\title{Mapeamento do Fluxo Ofensivo no Futebol:\\Uma Abordagem via PageRank}

% Remoção da enumeração e correção ortográfica (Allberti -> Alberti, Josue -> Josué)
\author[E. Alberti, J. Santos, L. Bicca \& L. Porto]{%
        Erick Alberti\footnote{\label{fgv}Fundação Getulio Vargas, Rio de Janeiro, Brasil},
        Josué Santos\textsuperscript{\ref{fgv}},
        Leonardo Bicca\textsuperscript{\ref{fgv}} \&
        Luigi Porto\textsuperscript{\ref{fgv}}}
\institute{Apresentado por: Erick Alberti, Josué Santos, Leonardo Bicca e Luigi Porto}
\date{19 de novembro de 2026}

\begin{document}

\begin{frame}[plain]
  \titlepage
\end{frame}

\section{Introdução e Motivação}

% =====================================================
\subsection{O Algoritmo PageRank}
\subsubsection{O que é o PageRank}
% =====================================================

\begin{frame}{O algoritmo PageRank}
\begin{itemize}
    \setlength{\itemsep}{\medskipamount} % Automatiza o espaçamento
    \item Desenvolvido para classificar páginas da internet
    segundo sua relevância, sendo um dos fundamentos
    originais do mecanismo de busca do Google.

    \item A estrutura da internet é representada por um
    grafo direcionado, no qual:
    \begin{itemize}
        \item Os vértices representam páginas;
        \item As arestas representam links entre páginas.
    \end{itemize}

    \item O princípio fundamental do algoritmo é que
    uma página é considerada mais relevante quando
    recebe referências de outras páginas relevantes.

    \item O resultado é uma distribuição de pontuações,
    denominada \textit{PageRank}, que permite
    ordenar os vértices segundo sua relevância.
\end{itemize}
\end{frame}

% =====================================================

% -----------------------------------------------------------
% SEÇÃO 2: MODELAGEM E FUNDAMENTAÇÃO MATEMÁTICA
% -----------------------------------------------------------
\section{Modelagem e Fundamentação Matemática}
  
  \subsection{Definições de Teoria dos Grafos}

    \begin{frame}{Estruturas Básicas de Teoria dos Grafos}
      Para compreender redes complexas, iniciamos com as definições estruturais mais fundamentais:
      \vfill
      \begin{itemize}
        \setlength{\itemsep}{\medskipamount}
        \item \textbf{Grafo ($G$):} Um par ordenado $G = (V, E)$, onde $V$ é um conjunto não-vazio de \textbf{vértices} (ou nós) e $E$ é um conjunto de \textbf{arestas} que conectam pares de vértices.
        \item \textbf{Digrafo (Grafo Direcionado):} Um grafo $D = (V, A)$ cujas conexões possuem direção. O conjunto $A$ é formado por \textbf{arcos} (arestas direcionadas), representados por pares ordenados $(u, v)$.
        \item \textbf{Grau de um Vértice:} O número de arestas incidentes a ele. Em digrafos, distingue-se o \textbf{grau de saída} ($d^+$) e o \textbf{grau de entrada} ($d^-$).
      \end{itemize}
    \end{frame}

    \begin{frame}{Representação Matricial e Conectividade}
      \begin{itemize}
        \setlength{\itemsep}{\medskipamount}
        \item \textbf{Matriz de Adjacência ($A$):} Uma matriz $n \times n$ (onde $n = |V|$) que representa as conexões do grafo de forma algébrica. Para um digrafo:
          \begin{equation*}
             A_{ij} = \begin{cases} 1, & \text{se existe um arco de } v_j \text{ para } v_i \\ 0, & \text{caso contrário} \end{cases}
          \end{equation*}
        \item \textbf{Conectividade Forte:} Um digrafo é fortemente conexo se, para todo par de vértices distintos $u, v \in V$, existe um caminho direcionado de $u$ para $v$ e um de $v$ para $u$. Isso assegura a ausência de nós sem saída na rede.
      \end{itemize}
    \end{frame}

  \subsection{Definições de Álgebra Linear}

    \begin{frame}{Vetores e Matrizes como Operadores}
      As definições estritamente necessárias para operar sobre a estrutura do grafo:
      \vfill
      \begin{itemize}
        \setlength{\itemsep}{\medskipamount}
        \item \textbf{Espaço Vetorial e Vetores:} Elementos de $\mathbb{R}^n$, representados como vetores-coluna. Neste contexto, atuam para denotar estados do sistema ou distribuições de probabilidade sobre os vértices.
        \item \textbf{Matriz Estocástica ($M$):} Matriz quadrada em que a soma dos elementos de cada coluna é exatamente 1 ($\sum_{i} M_{ij} = 1$) e onde $M_{ij} \ge 0$. Ela atua como o operador linear que mapeia o estado atual da rede para o próximo.
      \end{itemize}
    \end{frame}

    \begin{frame}{Autovalores, Autovetores e o Teorema Fundamental}
      
      \begin{definicao}[Autovalores e Autovetores]
        Um escalar $\lambda$ é um autovalor de uma matriz quadrada $M$ se existe um vetor não-nulo $x$ (o autovetor) tal que:
        \begin{equation}
            M x = \lambda x
        \end{equation}
      \end{definicao}
      
      \vfill
      
      \begin{teorema}[Perron-Frobenius]
        Para matrizes estocásticas por colunas com entradas estritamente positivas:
        \begin{itemize}
            \item O maior autovalor (autovalor dominante) é sempre $\lambda = 1$.
            \item Existe um autovetor único correspondente a $\lambda = 1$ cujas entradas são todas estritamente positivas.
        \end{itemize}
      \end{teorema}
      
    \end{frame}
\subsubsection{A estrutura matemática do PageRank}
% =====================================================

\begin{frame}{O modelo probabilístico}
\begin{itemize}
    \setlength{\itemsep}{\medskipamount}
    \item Considere um grafo direcionado com $n$ vértices.
    Seja $X_k$ a variável aleatória que representa
    o vértice ocupado após $k$ movimentos.

    \item A distribuição de probabilidade de $X_k$
    é representada pelo vetor:
    \[
    \mathbf{r}^{(k)} =
    \begin{pmatrix}
    \mathbb{P}(X_k=1)\\
    \mathbb{P}(X_k=2)\\
    \vdots\\
    \mathbb{P}(X_k=n)
    \end{pmatrix}
    \]

    \item Admitindo a propriedade de Markov, definimos
    a matriz de transição $P$, cujas entradas são:
    \[
    P_{ij} = \mathbb{P}(X_{k+1}=i\mid X_k=j)
    \]

    \item Assim, $P_{ij}$ representa a probabilidade
    de transição do vértice $j$ para o vértice $i$,
    com $\sum_{i=1}^{n}P_{ij}=1$.
\end{itemize}
\end{frame}

% =====================================================

\begin{frame}{A estrutura matemática do PageRank}
\begin{itemize}
    \setlength{\itemsep}{\medskipamount}
    \item O algoritmo introduz o teletransporte:
    com probabilidade $\alpha$, seguimos uma aresta;
    com probabilidade $1-\alpha$, escolhemos um vértice
    segundo a distribuição $\mathbf{v}$.

    \item A matriz de Google, que incorpora ambos
    os mecanismos, é definida por:
    \[
    G = \alpha P + (1-\alpha)\mathbf{v}\mathbf{1}^{T}
    \]

    \item O PageRank é a distribuição estacionária
    dessa matriz de transição, isto é, um vetor
    $\mathbf{r}$ que satisfaz:
    \[
    G\mathbf{r}=\mathbf{r}
    \]

    \item Portanto, $\mathbf{r}$ é um autovetor de $G$
    associado ao autovalor $1$, normalizado de modo que:
    \[
    \sum_{i=1}^{n}r_i=1
    \]
\end{itemize}
\end{frame}



\subsection{Dedução da fórmula do PageRank}

% =====================================================
\subsubsection{O caso mais simples}
% =====================================================

\begin{frame}{O caso mais simples: sem teletransporte}
\begin{itemize}
    \setlength{\itemsep}{\medskipamount}
    \item Queremos determinar a distribuição de probabilidade
    após $k+1$ movimentos, a partir da distribuição no instante $k$.

    \item Pela lei da probabilidade total:
    \[
    \mathbb{P}(X_{k+1}=i) = \sum_{j=1}^{n} \mathbb{P}(X_{k+1}=i\mid X_k=j) \mathbb{P}(X_k=j)
    \]

    \item Utilizando as definições anteriores:
    \[
    r_i^{(k+1)} = \sum_{j=1}^{n}P_{ij}r_j^{(k)}
    \]

    \item Portanto, em notação matricial:
    \[
    \boxed{ \mathbf{r}^{(k+1)}=P\mathbf{r}^{(k)} }
    \]
\end{itemize}
\end{frame}


% =====================================================
\subsubsection{O caso geral (com teletransporte)}
% =====================================================

\begin{frame}{O caso geral: com teletransporte}
\begin{itemize}
    \setlength{\itemsep}{\medskipamount}
    \item Introduzimos dois mecanismos de movimento:
    seguir uma aresta ($M$) ou realizar um teletransporte ($T$).

    \item Pela lei da probabilidade total:
    \[
    \begin{aligned}
    \mathbb{P}(X_{k+1}=i) ={}& \mathbb{P}(X_{k+1}=i\mid M)\alpha\\
    &+\mathbb{P}(X_{k+1}=i\mid T)(1-\alpha)
    \end{aligned}
    \]

    \item Substituindo as probabilidades de cada mecanismo:
    \[
    r_i^{(k+1)} = \alpha\sum_{j=1}^{n}P_{ij}r_j^{(k)} +(1-\alpha)v_i
    \]

    \item Finalmente, em notação matricial:
    \[
    \boxed{ \mathbf{r}^{(k+1)} = \alpha P\mathbf{r}^{(k)} +(1-\alpha)\mathbf{v} }
    \]
\end{itemize}
\end{frame}

% -----------------------------------------------------------
% SEÇÃO 4: APLICAÇÃO E ANÁLISE DE RESULTADOS
% -----------------------------------------------------------
\section{Aplicação e Análise de Resultados}

  \begin{frame}{Resultados do Mapeamento Ofensivo}
    % Substitua este slide pelas tabelas ou gráficos da sua análise de PageRank
    \begin{itemize}
        \item Espaço reservado para visualização do grafo de passes.
        \item Espaço reservado para a tabela dos jogadores com maior escore de PageRank no fluxo ofensivo.
    \end{itemize}
  \end{frame}

% -----------------------------------------------------------
% SEÇÃO 5: CONCLUSÃO E TRABALHOS FUTUROS
% -----------------------------------------------------------
\section{Conclusão e Trabalhos Futuros}

  \begin{frame}{Conclusões}
    % Substitua este slide pelas conclusões
    \begin{itemize}
        \item Espaço reservado para as considerações finais.
        \item Espaço reservado para possíveis melhorias e trabalhos futuros na modelagem.
    \end{itemize}
  \end{frame}

\section{Referências}

\begin{frame}[plain]
  \begin{fgvfocus}
    Obrigado! \\
    Perguntas?
  \end{fgvfocus}
\end{frame}

\end{document}